\section{CSV input and numeric table columns}
\label{sec:csv-input}
CSV files are read as characters, not executed as TeX. UTF-8 BOM, quoted commas,
doubled quotes and custom missing tokens are supported. Fields must fit on a
single line and the delimiter must be a comma. Headers must be unique and nonempty.
\subsection{Binding a native gradtable data source}
\label{sec:csv-native}
Load a dataset, then bind it with \key{data=results}. A column's
\key{source=Score} selects an exact CSV header and defaults to its id. Fields may
be reordered or reused. Omit \key{gradrow} in imported tables; rows follow file order,
numeric missing values use the dataset's mapping and text stays literal.
\begin{demo}{CSV field binding with multiple metrics}
\gradbarsloadcsv{results}{docs/data/results.csv}
\begin{gradtable}[data=results,width=\linewidth,scale=column,header align=c]
  \gradcolumn[source=Model]{method}{Method}
  \gradcolumn[type=bar,source=Score,mark=both]{score}{Score}
  \gradcolumn[type=number,source=Latency,
    options={precision=0}]{time}{Latency}
\end{gradtable}
\end{demo}
Manual \key{gradrow} data cannot be mixed with an imported dataset. Unknown datasets
or fields are errors. Quote list fields containing commas; use NA for missing
observations or endpoints within compound fields. Native sources support
\key{long=true}, with shared widths, scales, ranks and repeated headers across pages.
See Section~\ref{sec:column-scale} for automatic ranges and Section~\ref{sec:heat}
for changing a scalar column to \key{type=heat}.
\subsection{Sorting, filtering and displaying the first rows}
Sort names a declared column id; order selects ascending or descending. Numeric
sorting uses raw values, places missing last in either direction and retains
file order for ties. Filter combines comma-separated conditions with AND. Limit
is a nonnegative integer; 0 is unlimited. Processing order is filter, stable sort,
then limit.
\begin{demo}{Filter scores, then select the two lowest latencies}
\gradbarsloadcsv{results}{docs/data/results.csv}
\begin{gradtable}[data=results,width=\linewidth,scale=column,
  filter={score>=75},sort=time,order=ascending,limit=2,header align=c]
  \gradcolumn[source=Model]{method}{Method}
  \gradcolumn[type=bar,source=Score]{score}{Score}
  \gradcolumn[type=bar,source=Latency,
    options={better=lower,precision=0}]{time}{Latency}
\end{gradtable}
\end{demo}
Numeric conditions support =, !=, <, <=, > and >=; for example,
\key{filter={score>=75,time<100}}. Ordinary comparisons, including !=, exclude missing
values; use score=NA or score!=NA to test mapped missingness. Text conditions
support literal, case-sensitive = and !=; method= matches empty text. Protect a
whole condition containing a comma, as in \key{filter={{method=Compact, tuned}}}.
Text sorts in character order, not locale-specific collation, with empty text
last. These keys require data and text, number, bar, bullet or heat columns.
The loaded dataset remains unchanged; labels are not used to compare values.
\subsection{Scale scope and rankings after selection}
\key{scale scope=data} is the default: ranges still cover all imported rows.

Latency above uses [0,125]; \key{scale scope=visible} uses the retained two rows and
gives [0,88]. Rankings always use displayed rows. An empty visible selection uses
[0,1]. Long tables paginate after selection and retain shared ranges and rankings.


\subsection{Multiple keys and grouping order}
\begin{sloppypar}
\key{sort keys={score/descending,time/ascending}} sorts by score first and time
second. Each key requires its own direction. Nonempty sort and sort keys cannot
be combined; order applies only to the old single-key option. Missing stays last
at each key; equal/missing keys proceed to the next key and full ties keep CSV order.
\end{sloppypar}
\begin{demo}{Global top-three selection before grouping}
\gradbarsloadcsv{families}{docs/data/families.csv}
\begin{gradtable}[data=families,width=\linewidth,
  sort keys={score/descending,time/ascending},limit=3,group by=family,header align=c]
  \gradcolumn[source=Method]{method}{Method}
  \gradcolumn[source=Family]{family}{Family}
  \gradcolumn[type=number,source=Score]{score}{Score}
  \gradcolumn[type=number,source=Time]{time}{Time}
\end{gradtable}
\end{demo}
Selection runs filter, global stable sort, global limit, then grouping.
The selected rows are Compact, Wide and Fast; the original group order displays
Classic's Fast first, then Neural's Compact and Wide. Limit is a table-wide data
row count, not a per-group quota. Grouping changes whole-table order while
preserving relative sort order within each group.

\subsection{Legacy CSV access and numeric collection columns}

The following source assumes a CSV file with Name and Score columns:
\begin{source}
\gradbarsloadcsv{results}{results.csv}
\gradbarscsvtable{results}{Name}{Score}
\gradbarscsvstyle{shared}{results}{Score}
\gradbarcsv[style=shared]{results}{1}{Score}
\gradbarscsvcell{results}{1}{Name}
\end{source}
Rows are one-based. CSV range styles include zero and use a positive upper limit
even for all-zero, all-negative or all-missing columns. The legacy \key{gradbarscsvtable}
does not paginate; native \key{gradtable} data sources support \key{long=true}. A dataset definition follows grouping.

\begin{demo}{Only numbers in a common-style column}
\gradbarscolumn{G}{width=35mm,min=-50,max=100,precision=0}
\begin{tabular}{@{}lG@{}}
\toprule
\multicolumn{1}{c}{Item} & \multicolumn{1}{c}{Value}\\\midrule
A & 72\\ B & -20\\ C & 0\\ D & NA\\
\bottomrule
\end{tabular}
\end{demo}
Numeric columns use collcell. Use multicolumn for text headings. They do not
calculate a range and are not siunitx S columns. Avoid placing complete drawing
commands in a numeric collection cell. tabular and longtable workflows are
documented; a tabularray-specific collection interface is not implemented.
